Sei $m_\mathrm{G}$ die Masse des Gegengewichts, $c$ sein Abstand vom Turm, $m$ die Masse der Last und $r$ ihr Abstand vom Turm. Drehmoment des Gegengewichts = Drehmoment der Last:
\begin{align}
 m\,g\cdot r &= m_\mathrm{G}\,g\cdot c \quad\Rightarrow\quad m = \frac{m_\mathrm{G}\cdot c}{r}
\end{align}
Die Fallbeschleunigung kürzt sich weg.

\begin{abcliste}
\abc
\begin{align}
 m_1 &= \mLaF \\
   &= \frac{\mG\cdot\cG}{\ra} \\
   &= \mLa \approx \result{\mLatP{0}{2}}
\end{align}

\abc
\begin{align}
 m_2 &= \mLbF \\
   &= \frac{\mG\cdot\cG}{\rb} \\
   &= \mLb \approx \result{\mLbtP{0}{3}}
\end{align}
Darum erlaubt die Lasttabelle eines Krans umso weniger, je weiter draussen die Last hängt.
\end{abcliste}
